\documentclass[12pt]{article}
\usepackage[utf8]{inputenc}
\usepackage[T1]{fontenc}
\usepackage{hyperref}
\title{Kabix\'i Language Resources for K-12 Curriculum Development}
\author{Compiled by Claude Sonnet 5.5 (claude-sonnet-5-5) for Andr\'as Kornai}
\date{}
\begin{document}
\maketitle

\section{Introduction}
Kabix\'i (ISO 639-3 xbx; Glottocode kabi1264) is a poorly documented South American name. The search results do not agree on which language it refers to. Wikipedia describes an extinct Tupian language called Kabixiana in Rondônia, and a Chapacuran language called Wanyam, whose dialects include Cabixi. It also describes a separate unclassified Cabixi language of Mato Grosso, known only from a word list collected by Natterer. One Wikipedia-derived summary says the code xbx is retired. These identifications were not reconciled, and I make no claim about which is the intended one. No living speakers, curricula or teaching materials were found.

\section{Primary Educational Resources}
None found.

\section{Linguistic and Reference Materials}
\begin{itemize}
\item Kabixiana language (Wikipedia): \url{https://en.wikipedia.org/wiki/Kabixiana_language}
\item Cabixi language (Wikipedia), Mato Grosso: \url{https://en.wikipedia.org/wiki/Cabixi_language}
\item Wanyam language (Wikipedia), Chapacuran: \url{https://en.wikipedia.org/wiki/Wanyam_language}
\end{itemize}

\section{Religious and Community Resources}
Joshua Project has a Kabixi people profile: \url{https://joshuaproject.net/people_groups/19322}. Its classification conflicts with the Wikipedia entries, so treat it with caution. No scripture or community language organisation was found.

\section{Digital and Audio Resources}
None found.

\section{Curriculum Development Context}
Only short word lists appear to exist, and the identity of the language is unclear. A curriculum is not currently feasible beyond a cultural-history unit. Further work should start by establishing which language and community are meant, ideally with Chapacuran-speaking communities in Rondônia and Mato Grosso, and with specialist linguists.
\end{document}